\glsresetall

\chapter{Evaluation}
This chapter evaluates the effort required to extend the Katamaran \gls{pmp} case study with support for \gls{smepmp}.
The evaluation is based on a comparison between the original case study and the extended implementation, allowing the scope of the required modifications to be quantified.

The changes are presented per \gls{smepmp} feature, following the structure introduced in the previous chapter.
Finally, the chapter discusses the experience of extending the Katamaran framework, including observations on its usability and the challenges encountered during development.

\section{Modification Effort}
To extend the existing Katamaran \gls{pmp} case study with support for \gls{smepmp}, five source files required modification.
In total, 112 lines of code were added, 14 lines were removed, and one existing line was modified.
\Cref{table:results} provides a breakdown of these changes by feature.

The table only includes modifications made to the implementation and the \texttt{mseccfg} register contracts.
The contracts corresponding to the semantic changes introduced by \gls{mmwp}, \gls{mml} and \gls{rlb} were not updated as part of this work and are thus not included in the reported figures.

\begin{table}[]
\centering
\begin{tabular}{c|ccc}
                            & Added & Removed & Changed \\
\hline
\texttt{mseccfg}            & 28    & 0       & 0       \\
\texttt{mseccfg} contracts  & 47    & 0       & 1       \\
\gls{mmwp}                  & 13    & 13      & 0       \\
\gls{mml}                   & 58    & 0       & 0       \\
\gls{rlb}                   & 14    & 1       & 0
\end{tabular}
\caption{Updates to Katamaran}
\label{table:results}
\end{table}

As shown in \cref{table:results}, the required effort differs considerably between the individual \gls{smepmp} features.
While support for \gls{rlb} required only a small number of changes, \gls{mml} required substantially more modifications because it fundamentally changes the interpretation of existing \gls{pmp} permissions.
The following sections discuss the changes required for each feature in more detail.

\subsection{\texttt{mseccfg} Control Status Register}
As shown in \cref{table:results}, adding support for the mseccfg \gls{csr} required 28 lines of implementation code.
Most of these additions consist of defining the new register, its layout, and the associated field definitions used throughout the model.

The corresponding contract updates are more extensive, requiring 47 additional lines.
This is because multiple functions interact with the \glspl{csr}, and their contracts must be updated to include the mseccfg register in the formal specification.
Although the implementation changes are localized, the new architectural state must be reflected consistently in every affected contract.

\subsection{Machine-Mode Whitelist Policy}
As shown in \cref{table:results}, adding support for \gls{mmwp} required 13 additions and 13 removals.
The identical number of additions and removals is coincidental and results from replacing existing privilege-level checks with updated checks that also account for the state of \gls{mmwp}.

Before the introduction of \gls{mmwp}, the model allowed Machine-mode accesses when no \gls{pmp} region matched, while accesses from lower privilege modes were denied.
Supporting \gls{mmwp} required modifying these checks such that unmatched accesses from Machine-mode are only allowed when \gls{mmwp} is disabled.
Although the number of changed lines is relatively small, the modification affects a central part of the access-control logic because it changes the default behavior for Machine-mode memory accesses.

\subsection{Machine-Mode Lockdown}
As shown in \cref{table:results}, adding support for \gls{mml} required 58 additions, making it the largest individual change among the \gls{smepmp} features.
This is because \gls{mml} fundamentally changes the interpretation of existing \gls{pmp} entries rather than only adding an additional condition to the access checks.

Before the introduction of \gls{mml}, the model determined access permissions based on the privilege mode and the R-, W-, and X-bits from the corresponding \texttt{pmpcfg} register.
With \gls{mml} enabled, the interpretation of the L-bit changes and additional cases must be handled to support Machine-mode-only, User- and Supervisor-mode-only, and shared memory regions.
As a result, the permission evaluation logic requires additional checks before the existing R-, W-, and X-bit based permission checks can be applied.

\subsection{Rule Locking Bypass}
As shown in \cref{table:results}, supporting \gls{rlb} required 14 additions and 1 removal.
Compared to the other \gls{smepmp} features, this is a relatively small change because \gls{rlb} only affects how the model determines whether a \gls{pmp} entry is considered locked.
Furthermore, the original RISC-V model and thus the Katamaran model already encapsulate this check in a dedicated helper function.
As a result, supporting \gls{rlb} only required modifying this function, without affecting the remainder of the permission evaluation logic.

Prior to the introduction of \gls{rlb}, whether a \gls{pmp} entry was locked was determined solely by the L-bit in the corresponding \texttt{pmpcfg} register.
To support \gls{rlb}, this check was extended such that locked entries are treated as unlocked when the \gls{rlb} bit is enabled.

% TODO add more observations? (don't know what)
\section{Discussion on Katamaran use}
Extending the Katamaran \gls{pmp} case study with support for \gls{smepmp} provided several insights into the development experience of working with the framework.

Before modifying the existing case study, time was spent learning the basics of Rocq, the proof assistant on which Katamaran is built.
This provided a useful introduction to the proof language and tactics used throughout the framework.
However, familiarity with Rocq alone was not sufficient to understand the Katamaran case study, as most of the implementation is written in \textmu{}Sail.
Consequently, understanding the structure and semantics of \textmu{}Sail proved to be equally important.

Since the \textmu{}Sail model is derived from the official RISC-V Sail model, consulting the original Sail implementation was often helpful when interpreting the existing code.
In many cases, the correspondence between the two models was straightforward.
However, some parts of the \textmu{}Sail implementation differed sufficiently from the original Sail code that identifying the corresponding functionality required additional investigation, increasing the time needed to understand and modify the implementation.
